% Фиксированные схемы урока 3, скопированные из задач Ященко-2022.
% В аргументах макросов допускается менять только числовые подписи.

% Варианты 1--2: тупоугольный равнобедренный треугольник с высотой.
\newcommand{\figIsoscelesAltitude}{%
\begin{tikzpicture}[scale=1.15,line cap=round,line join=round,every node/.style={font=\small}]
  \coordinate (A) at (0,0); \coordinate (B) at (5.2,0); \coordinate (C) at (2.9,1.25);
  \coordinate (H) at ($(B)!(A)!(C)$);
  \draw[line width=0.9pt] (A)--(B);
  \draw[line width=0.9pt,equal segment] (A)--(C);
  \draw[line width=0.9pt,equal segment] (C)--(B);
  \draw[dashed,line width=0.8pt] (A)--(H);
  \draw[dashed,line width=0.8pt] (H)--(C);
  \pic[draw,line width=0.7pt,angle radius=3.5mm] {right angle=A--H--C};
  \node[below left] at (A) {$A$}; \node[below right] at (B) {$B$};
  \node[above right] at (C) {$C$}; \node[above left] at (H) {$H$};
\end{tikzpicture}}

% Вариант 7: высота и медиана из вершины прямого угла.
\newcommand{\figRightAltitudeMedian}[1]{%
\begin{tikzpicture}[scale=0.9,line cap=round,line join=round,every node/.style={font=\small}]
  \coordinate (A) at (0,0); \coordinate (B) at (6,0); \coordinate (M) at (3,0);
  \coordinate (H) at (3.521,0); \coordinate (C) at (3.521,2.954);
  \draw[line width=0.9pt,equal segment] (A)--(M);
  \draw[line width=0.9pt,equal segment] (M)--(B);
  \draw[line width=0.9pt] (A)--(C)--(B);
  \draw[line width=0.85pt] (C)--(M); \draw[line width=0.85pt] (C)--(H);
  \pic[draw,line width=0.65pt,angle radius=3.5mm] {right angle=A--C--B};
  \pic[draw,line width=0.65pt,angle radius=3.2mm] {right angle=C--H--B};
  \pic[draw,line width=0.65pt,angle radius=7mm] {angle=C--B--A};
  \node at (5.15,0.48) {$#1^\circ$};
  \node[below left] at (A) {$A$}; \node[below right] at (B) {$B$};
  \node[above] at (C) {$C$}; \node[below] at (M) {$M$}; \node[below] at (H) {$H$};
\end{tikzpicture}}

% Вариант 8: биссектриса и медиана из вершины прямого угла.
\newcommand{\figRightBisectorMedian}[1]{%
\begin{tikzpicture}[scale=0.9,line cap=round,line join=round,every node/.style={font=\small}]
  \coordinate (A) at (0,0); \coordinate (B) at (6,0); \coordinate (M) at (3,0);
  \coordinate (D) at (3.528,0); \coordinate (C) at (4.026,2.819);
  \draw[line width=0.9pt,equal segment] (A)--(M);
  \draw[line width=0.9pt,equal segment] (M)--(B);
  \draw[line width=0.9pt] (A)--(C)--(B);
  \draw[line width=0.85pt] (C)--(M); \draw[line width=0.85pt] (C)--(D);
  \pic[draw,line width=0.65pt,angle radius=3.5mm] {right angle=A--C--B};
  \pic[draw,line width=0.65pt,angle radius=7mm] {angle=A--C--D};
  \pic[draw,line width=0.65pt,angle radius=7mm] {angle=D--C--B};
  \pic[draw,line width=0.65pt,angle radius=4.5mm] {angle=M--C--D};
  \node at (3.68,2.10) {$#1^\circ$};
  \node[below left] at (A) {$A$}; \node[below right] at (B) {$B$};
  \node[above] at (C) {$C$}; \node[below] at (M) {$M$}; \node[below] at (D) {$D$};
\end{tikzpicture}}

% Вариант 12: равнобедренный треугольник с высотой.
\newcommand{\figIsoscelesSine}[2]{%
\begin{tikzpicture}[scale=0.72,line cap=round,line join=round,every node/.style={font=\small}]
  \coordinate (A) at (0,0); \coordinate (H) at (5.292,0); \coordinate (B) at (6.047,0); \coordinate (C) at (5.292,6);
  \draw[line width=0.9pt,equal segment] (A)--(B);
  \draw[line width=0.9pt,equal segment] (B)--(C);
  \draw[line width=0.9pt] (C)--(A);
  \draw[line width=0.85pt] (C)--(H);
  \pic[draw,line width=0.65pt,angle radius=3.2mm] {right angle=C--H--B};
  \path (C)--(H) node[midway,right=-5mm] {$#1$};
  \path (A)--(C) node[midway,above left] {$#2$};
  \node[below left] at (A) {$A$}; \node[below right] at (B) {$B$};
  \node[above] at (C) {$C$}; \node[below] at (H) {$H$};
\end{tikzpicture}}

% Вариант 11: биссектрисы в треугольнике.
\newcommand{\figIncenterBisectors}[1]{%
\begin{tikzpicture}[scale=0.75,line cap=round,line join=round,every node/.style={font=\small}]
  \coordinate (A) at (0,0); \coordinate (B) at (8,0); \coordinate (C) at (3.5,9);
  \coordinate (D) at (5.96,4.08); \coordinate (E) at (1.55,3.99); \coordinate (O) at (3.80,2.60);
  \draw[line width=0.9pt] (A)--(B)--(C)--cycle;
  \draw[line width=0.85pt] (A)--(D); \draw[line width=0.85pt] (B)--(E);
  \fill (O) circle (1.4pt);
  \pic[draw,line width=0.6pt,angle radius=6mm] {angle=B--A--D};
  \pic[draw,line width=0.6pt,angle radius=6mm] {angle=D--A--C};
  \pic[draw,line width=0.6pt,angle radius=5mm] {angle=C--B--E};
  \pic[draw,line width=0.6pt,angle radius=6.4mm] {angle=C--B--E};
  \pic[draw,line width=0.6pt,angle radius=5mm] {angle=E--B--A};
  \pic[draw,line width=0.6pt,angle radius=6.4mm] {angle=E--B--A};
  \pic[draw,line width=0.65pt,angle radius=7mm] {angle=A--C--B};
  \node at (3.55,7.15) {$#1^\circ$};
  \node[below left] at (A) {$A$}; \node[below right] at (B) {$B$}; \node[above] at (C) {$C$};
  \node[right] at (D) {$D$}; \node[left] at (E) {$E$}; \node[below] at (O) {$O$};
\end{tikzpicture}}

% Вариант 25: высота и биссектриса из вершины прямого угла.
\newcommand{\figRightAltitudeBisector}{%
\begin{tikzpicture}[scale=0.95,line cap=round,line join=round,every node/.style={font=\small}]
  \coordinate (A) at (0,0); \coordinate (B) at (6,0); \coordinate (C) at (5.007,2.229);
  \coordinate (H) at (5.007,0); \coordinate (D) at (4.152,0);
  \draw[line width=0.9pt] (A)--(B)--(C)--cycle;
  \draw[line width=0.85pt] (C)--(H); \draw[line width=0.85pt] (C)--(D);
  \pic[draw,line width=0.65pt,angle radius=3.3mm] {right angle=A--C--B};
  \pic[draw,line width=0.65pt,angle radius=3mm] {right angle=C--H--B};
  \node[below left] at (A) {$A$}; \node[below right] at (B) {$B$};
  \node[above] at (C) {$C$}; \node[below] at (D) {$D$}; \node[below] at (H) {$H$};
\end{tikzpicture}}

% Вариант 32: две высоты в остроугольном треугольнике.
\newcommand{\figTriangleAltitudes}{%
\begin{tikzpicture}[scale=1.02,line cap=round,line join=round,every node/.style={font=\small}]
  \coordinate (A) at (0,0); \coordinate (B) at (7,0); \coordinate (C) at (1.544,3.821);
  \coordinate (D) at (2.303,3.289); \coordinate (E) at (0.982,2.431); \coordinate (H) at (1.544,2.205);
  \draw[line width=0.9pt] (A)--(B)--(C)--cycle;
  \draw[line width=0.85pt] (A)--(D); \draw[line width=0.85pt] (B)--(E);
  \pic[draw,line width=0.6pt,angle radius=3mm] {right angle=A--D--C};
  \pic[draw,line width=0.6pt,angle radius=3mm] {right angle=B--E--C};
  \pic[draw,line width=0.65pt,angle radius=7mm] {angle=B--A--C};
  \pic[draw,line width=0.65pt,angle radius=7mm] {angle=C--B--A};
  \pic[draw,line width=0.7pt,angle radius=5.5mm] {angle=A--H--B};
  \fill (H) circle (1.3pt);
  \node[below left] at (A) {$A$}; \node[below right] at (B) {$B$}; \node[above] at (C) {$C$};
  \node[above right] at (D) {$D$}; \node[left] at (E) {$E$}; \node[below] at (H) {$H$};
\end{tikzpicture}}

% Вариант 18: сторона треугольника в описанной окружности.
\newcommand{\figCircumcircleSide}[2]{%
\begin{tikzpicture}[scale=1.15,line cap=round,line join=round,every node/.style={font=\small}]
  \def\R{2.3}
  \coordinate (O) at (0,0); \coordinate (A) at (150:\R); \coordinate (B) at (30:\R); \coordinate (C) at (90:\R);
  \draw[line width=0.9pt] (O) circle (\R); \draw[line width=0.9pt] (A)--(B)--(C)--cycle;
  \pic[draw,line width=0.65pt,angle radius=5mm] {angle=A--C--B};
  \node at (0,1.40) {$#2^\circ$};
  \node[left] at (A) {$A$}; \node[right] at (B) {$B$}; \node[above] at (C) {$C$};
  \node[below] at (0,-2.45) {$R=#1$};
\end{tikzpicture}}
